\section{The construction in ZFC}
\label{sec:construction}

We now choose the bridge steps so that every pure state eventually
joins the class of the state peaked by $(p_n)$. Countable
determination ensures that every pure state of the final algebra
will have appeared early enough to be treated.

\begin{theorem}[Kishimoto--Ozawa--Sakai]
\label{na:homogeneity}
Let $B$ be a separable $C^*$-algebra. If pure states $h,k$ of $B$
have GNS representations with the same kernel, there is an
automorphism $\gamma$ of $B$ such that $h\circ\gamma=k$.
In particular, this holds for every pair of pure states of a
separable simple $C^*$-algebra.
\end{theorem}

This is \cite[Theorem~1.1]{KishimotoOzawaSakai2003}, which in fact
provides an asymptotically inner automorphism. We apply it only to
the separable algebras $A[S]$ of \cref{na:inherited}.

\begin{proof}[Proof of \cref{na:main}]
Set $\kappa=\mathfrak c^+$. Every ordinal below $\kappa$ has
cardinality at most $\mathfrak c$, and every countable subset of
$\kappa$ is bounded: otherwise $\kappa$ would be a countable union
of sets of cardinality at most $\mathfrak c$, a contradiction.

\smallskip\noindent
\emph{Fix the schedule first.}
Before constructing any algebras, choose an injection
\begin{equation}\label{eq:schedule}
 \sigma:\kappa\times\kappa\longrightarrow\kappa,
 \qquad \sigma(\beta,\eta)\ge\beta.
\end{equation}
To do so, enumerate $\kappa\times\kappa$ in order type $\kappa$.
When a pair $(\beta,\eta)$ is encountered, choose an unused ordinal
at least $\beta$. Fewer than $\kappa$ ordinals have been used,
whereas the tail above $\beta$ has cardinality $\kappa$.
This defines the entire schedule without referring to any states.

\smallskip\noindent
\emph{Construct the algebras in ordinal order.}
Build a continuous bridge iteration $(X_\beta)_{\beta\le\kappa}$
starting from the algebra of \cref{na:base-peak}. For each
$\beta<\kappa$, once $X_\beta$ has been constructed, choose a list
\[
 (h_{\beta,\eta})_{\eta<\kappa}
\]
containing every pure state of $X_\beta$, with repetitions allowed.
This is possible by \cref{na:state-count}, since
$|\beta|\le\mathfrak c$ for $\beta<\kappa$.

At stage $j<\kappa$, let $f$ be the state of $X_j$ peaked by
$(p_n)$. If $j$ is outside the range of $\sigma$, make the identity
step. Otherwise write $j=\sigma(\beta,\eta)$ and let $g$ be the
unique lift of $h_{\beta,\eta}$ to $X_j$. The list for $\beta$ is
already available because $\beta\le j$; when $\beta=j$, it is
chosen before the successor step. If $g\sim f$, again make the
identity step.

Suppose $g\not\sim f$. By \cref{na:countable-determination}, choose
a countable dependency-closed $S\subset j$ such that $g|_{A[S]}$
is pure and uniquely determines $g$ on $X_j$. The restriction
$f|_{A[S]}$ is pure by \cref{na:peaking-restriction}. The algebra
$A[S]$ is separable, simple and unital by \cref{na:inherited}.
Applying \cref{na:homogeneity}, choose an automorphism $\gamma$
of $A[S]$ with
\[
 (g|_{A[S]})\circ\gamma=f|_{A[S]},
 \qquad q_n^j=\gamma(p_n)\quad(n\ge0).
\]
If a state on $A[S]$ takes value one on every $q_n^j$, composing
it with $\gamma$ gives the unique state taking value one on every
$p_n$. Thus $(q_n^j)$ peaks to $g|_{A[S]}$. By
\cref{na:closed-extension}, it peaks to $g$ on $X_j$ as well.

The trace requirement deserves a separate check. For each $n$,
$\gamma$ carries the unital matrix stage $M_{2^n}\subset A_0$
onto a unital matrix subalgebra of $A[S]$, and carries its minimal
projection $p_n$ to $q_n^j$. The ambient trace restricts to the
unique normalized trace on this matrix algebra. Consequently,
\[
 \tau(q_n^j)=2^{-n}=\tau(p_n).
\]
This argument does not require an automorphism to preserve the
specified trace on all of $A[S]$.

We can therefore perform the bridge of \cref{na:bridge} with
$(p_n),(q_n^j)$, and set $I_j=S$. At an identity step set
$I_j=\varnothing$. At limit stages take the norm completion of
the union. Each successor satisfies \textup{(U)} and has a faithful
trace-preserving expectation onto its base. The results of
\cref{sec:iteration} therefore apply throughout the recursion.

All choices are set-sized. At every stage the algebra is generated
by at most $\kappa$ elements, so its density is at most $\kappa$
and its underlying cardinality at most $\lambda=\kappa^{\aleph_0}$.
Fix a set of cardinality $\lambda$. Algebra structures on its
subsets, together with their norms, traces, embeddings, states,
automorphisms and lists, form a set of codes. Well-order that set
to make the required choices; recode successor algebras with their
embeddings and use norm inductive limits. Ordinary choice and
transfinite recursion in ZFC suffice.

\smallskip\noindent
\emph{Every final pure state is treated.}
Let $A=X_\kappa$ and let $h$ be any pure state of $A$.
By \cref{na:countable-determination}, there is a countable
dependency-closed $S\subset\kappa$ such that $h|_{A[S]}$ is pure
and has unique extension $h$ to $A$. Boundedness of $S$ gives
$\beta<\kappa$ with $A[S]\subset X_\beta$.
The state $h|_{X_\beta}$ is the unique extension of $h|_{A[S]}$
to $X_\beta$: any other extension could itself be extended to
$A$, contradicting uniqueness there. Hence $h|_{X_\beta}$ is pure.

Choose $\eta$ with $h_{\beta,\eta}=h|_{X_\beta}$ and put
$j=\sigma(\beta,\eta)$. Its lift to $X_j$ is $h|_{X_j}$.
At stage $j$ this lift either already belongs to the $p_n$-class
or is joined to it by the bridge. By \cref{na:unique-chain}, the
equivalence persists to $A$. Thus all pure states of $A$ have
unitarily equivalent GNS representations.

Every nonzero irreducible representation $\pi$ of the unital
algebra $A$ is unital: the range of $\pi(1)$ is a nonzero closed
invariant subspace. Every nonzero vector is cyclic by irreducibility,
so each unit vector defines a pure state whose GNS representation
is equivalent to $\pi$. The $p_n$-state supplies one such
representation. Therefore $A$ has exactly one nonzero irreducible
representation up to unitary equivalence.

Finally, $A$ is simple and has a faithful tracial state by
\cref{na:inherited} and the limit permanence facts. It is unital
and contains the infinite-dimensional algebra $A_0$. An algebra
$K(H)$ is unital only when $H$ is finite-dimensional:
the identity on an infinite-dimensional Hilbert space is not
compact, as an infinite orthonormal sequence has no norm-convergent
subsequence. The finite-dimensional case, including $H=0$, cannot
be isomorphic to $A$. This proves the theorem.
\end{proof}
